\chapter{Midpoints and Parallels}\label{ch:g8:midpoints}

Join the midpoints of two sides of a triangle: the \omterm{def:g6:lines:objects}{segment} you get is
always \omterm{def:g6:lines:perp}{parallel} to the third side, and exactly \omterm{def:g3:division:half}{half} as long. This
``midpoint theorem'' and its converse are the first taste of a great
idea --- \omterm{def:g6:lines:perp}{parallels} cut lengths proportionally --- which blossoms into
Thales' theorem in \cref{ch:g9:thales}.

\section{The midpoint theorem}

\begin{theorem}[Midpoint theorem]\label{thm:g8:midpoints:direct}
In a triangle $ABC$, let $I$ be the midpoint of $[AB]$ and $J$ the
midpoint of $[AC]$. Then the line $(IJ)$ is \omterm{def:g6:lines:perp}{parallel} to $(BC)$, and
\[
IJ = \frac{BC}{2} .
\]
\end{theorem}

\begin{proof}[Idea of proof]
Let $K$ be the \omterm{def:g7:central:def}{symmetric} of $I$ about $J$ (a half-turn,
\cref{ch:g7:central}). The \omterm{def:g4:geometry:polygon}{quadrilateral} $AICK$ has diagonals $[IK]$
and $[AC]$ crossing at their common midpoint $J$: it is a
\omterm{def:g7:central:parallelogram}{parallelogram}, so $KC$ is \omterm{def:g6:lines:perp}{parallel} to $AI$, i.e.\ to $(AB)$, and
$KC = AI = IB$. Then $IBCK$ has two opposite sides ($[IB]$ and
$[KC]$) \omterm{def:g6:lines:perp}{parallel} and equal: it is a \omterm{def:g7:central:parallelogram}{parallelogram} too
(\cref{thm:g7:central:recognize}), so $(IK) \parallel (BC)$ ---
that is $(IJ) \parallel (BC)$ --- and $BC = IK = 2\,IJ$.
\end{proof}

\begin{omfigure}
\begin{tikzpicture}[scale=1.1]
  \coordinate (A) at (1.7,2.9);
  \coordinate (B) at (0,0);
  \coordinate (C) at (4.6,0);
  \coordinate (I) at ($(A)!0.5!(B)$);
  \coordinate (J) at ($(A)!0.5!(C)$);
  \draw[thick] (A) -- (B) -- (C) -- cycle;
  \draw[omThm, very thick] (I) -- (J);
  \fill (A) circle (1.5pt) node[above, font=\small] {$A$};
  \fill (B) circle (1.5pt) node[below left, font=\small] {$B$};
  \fill (C) circle (1.5pt) node[below right, font=\small] {$C$};
  \fill[omThm] (I) circle (1.5pt) node[left, font=\small] {$I$};
  \fill[omThm] (J) circle (1.5pt) node[right, font=\small] {$J$};
  \draw ($(A)!0.22!(B)$) ++(-0.09,-0.05) -- ++(0.18,0.1);
  \draw ($(A)!0.72!(B)$) ++(-0.09,-0.05) -- ++(0.18,0.1);
  \draw ($(A)!0.25!(C)$) ++(-0.06,-0.1) -- ++(0.12,0.2);
  \draw ($(A)!0.25!(C)$) ++(0.02,-0.1) -- ++(0.12,0.2);
  \draw ($(A)!0.75!(C)$) ++(-0.06,-0.1) -- ++(0.12,0.2);
  \draw ($(A)!0.75!(C)$) ++(0.02,-0.1) -- ++(0.12,0.2);
\end{tikzpicture}

{\small The midpoint \omterm{def:g6:lines:objects}{segment} $[IJ]$ (red) is \omterm{def:g6:lines:perp}{parallel} to the third
side $(BC)$ and \omterm{def:g3:division:half}{half} as long.}
\end{omfigure}

\begin{example}\label{ex:g8:midpoints:direct}
In a triangle $ABC$ with $BC = 9$ cm, the midpoints $I$ of $[AB]$ and
$J$ of $[AC]$ are joined. Without any measurement:
$(IJ) \parallel (BC)$ and $IJ = \frac92 = 4.5$ cm. The small triangle
$AIJ$ is a half-size copy of $ABC$ --- its \omterm{def:g6:measure:perimeter}{perimeter} is \omterm{def:g3:division:half}{half} the
\omterm{def:g6:measure:perimeter}{perimeter} of $ABC$ as well.
\end{example}

\begin{theorem}[Converse]\label{thm:g8:midpoints:converse}
In a triangle $ABC$, the line through the midpoint $I$ of $[AB]$ and
\emph{\omterm{def:g6:lines:perp}{parallel}} to $(BC)$ crosses the side $[AC]$ at its midpoint.
\end{theorem}
\admitted

\begin{example}\label{ex:g8:midpoints:converse}
A path crosses a triangular field, starting at the middle of one edge
and running \omterm{def:g6:lines:perp}{parallel} to the opposite edge. The converse guarantees
that the path exits exactly at the middle of the other edge --- no
measuring needed on the far side.
\end{example}

\begin{method}[Choosing between theorem and converse]\label{met:g8:midpoints:choose}
\begin{enumerate}
  \item \emph{Two midpoints known} $\Rightarrow$ use the direct
        theorem: conclude parallelism and half-length.
  \item \emph{One midpoint and a \omterm{def:g6:lines:perp}{parallel} known} $\Rightarrow$ use
        the converse: conclude that the crossing point is a midpoint.
  \item Write which triangle, which midpoints and which theorem you
        use --- one sentence each.
\end{enumerate}
\end{method}

\section{Half-size triangles}

\begin{proposition}[The midpoint triangle]\label{prop:g8:midpoints:medial}
Joining the three midpoints of the sides of a triangle cuts it into
four triangles of equal \omterm{def:g6:measure:area}{areas}, each a half-size copy of the original.
\end{proposition}

\begin{proof}
Each side of the midpoint triangle is \omterm{def:g6:lines:perp}{parallel} to a side of $ABC$ and
\omterm{def:g3:division:half}{half} as long (\cref{thm:g8:midpoints:direct}, applied three times).
The four small triangles have sides of the same three lengths, so
they are identical copies; together they tile $ABC$, so each has a
\omterm{def:g3:division:half}{quarter} of its \omterm{def:g6:measure:area}{area}.
\end{proof}

\begin{omfigure}
\begin{tikzpicture}[scale=1.05]
  \coordinate (A) at (1.9,3.0);
  \coordinate (B) at (0,0);
  \coordinate (C) at (5.0,0);
  \coordinate (I) at ($(A)!0.5!(B)$);
  \coordinate (J) at ($(A)!0.5!(C)$);
  \coordinate (K) at ($(B)!0.5!(C)$);
  \draw[thick] (A) -- (B) -- (C) -- cycle;
  \fill[omThm!12] (I) -- (J) -- (K) -- cycle;
  \draw[omThm, thick] (I) -- (J) -- (K) -- cycle;
  \fill (A) circle (1.4pt) node[above, font=\small] {$A$};
  \fill (B) circle (1.4pt) node[below left, font=\small] {$B$};
  \fill (C) circle (1.4pt) node[below right, font=\small] {$C$};
  \fill[omThm] (I) circle (1.4pt) node[left, font=\small] {$I$};
  \fill[omThm] (J) circle (1.4pt) node[right, font=\small] {$J$};
  \fill[omThm] (K) circle (1.4pt) node[below, font=\small] {$K$};
\end{tikzpicture}

{\small The midpoint triangle $IJK$ cuts $ABC$ into four equal
triangles --- a picture worth remembering.}
\end{omfigure}

\begin{remark}[Towards Thales]
The midpoint theorem is the case ``one \omterm{def:g3:division:half}{half}'' of a more general fact:
a line \omterm{def:g6:lines:perp}{parallel} to one side of a triangle cuts the two other sides in
\emph{equal ratios} --- one third and one third, two fifths and two
fifths\,\dots\ That general statement is Thales' theorem
(\cref{ch:g9:thales}); the midpoint case is the only one needed this
year.
\end{remark}

\section{Exercises}

\begin{exercise}[$\star$]\label{exo:g8:midpoints:1}
In a triangle $RST$, $M$ is the midpoint of $[RS]$ and $N$ the
midpoint of $[RT]$, with $ST = 7$ cm. What can be said of the line
$(MN)$ and the length $MN$? Cite the theorem used.
\end{exercise}

\begin{exercise}[$\star$]\label{exo:g8:midpoints:2}
In a triangle $ABC$, $I$ is the midpoint of $[AB]$ and the line
through $I$ \omterm{def:g6:lines:perp}{parallel} to $(BC)$ cuts $[AC]$ at $J$, with
$AC = 11$ cm. Compute $AJ$, citing the theorem used.
\end{exercise}

\begin{exercise}[$\star$]\label{exo:g8:midpoints:3}
Draw a triangle $ABC$ with $AB = 8$ cm, $AC = 6$ cm, $BC = 7$ cm,
place the midpoints $I$ of $[AB]$ and $J$ of $[AC]$, and measure $IJ$
to check the theorem. Compute the \omterm{def:g6:measure:perimeter}{perimeter} of $AIJ$ without
measuring anything else.
\end{exercise}

\begin{exercise}[$\star$]\label{exo:g8:midpoints:4}
$IJK$ is the midpoint triangle of $ABC$, whose sides measure $10$,
$12$ and $16$ cm. Give the three sides of $IJK$ and its \omterm{def:g6:measure:perimeter}{perimeter}.
How do the two \omterm{def:g6:measure:perimeter}{perimeters} compare?
\end{exercise}

\begin{exercise}[$\star$]\label{exo:g8:midpoints:5}
The \omterm{def:g6:measure:area}{area} of a triangle is $36$ cm$^2$. What is the \omterm{def:g6:measure:area}{area} of its
midpoint triangle? Of each of the four small triangles?
\end{exercise}

\begin{exercise}[$\star\star$]\label{exo:g8:midpoints:6}
In a triangle $ABC$, $I$ is the midpoint of $[AB]$, $J$ that of
$[AC]$ and $K$ that of $[BC]$. Show that $IJKB$ is a \omterm{def:g7:central:parallelogram}{parallelogram}.
(Compare $[IJ]$ and $[BK]$: \omterm{def:g6:lines:perp}{parallel}? equal?)
\end{exercise}

\begin{exercise}[$\star\star$]\label{exo:g8:midpoints:7}
A triangular sail has its lowest edge $6.4$ m long. A reinforcement
seam joins the midpoints of the two other edges. What length of seam
is needed? What if the seam instead joins the points located at one
\omterm{def:g3:division:half}{quarter} of each edge, starting from the top \omterm{def:g5:solids:def}{vertex}? (Conjecture with
a drawing; the proof is Thales', \cref{ch:g9:thales}.)
\end{exercise}

\begin{exercise}[$\star\star$]\label{exo:g8:midpoints:8}
In triangle $ABC$, \omterm{def:g6:shapes:triangles}{right-angled} at $A$, $M$ is the midpoint of the
hypotenuse $[BC]$, and the \omterm{def:g6:lines:perp}{parallel} to $(AB)$ through $M$ meets
$[AC]$ at $N$.
\begin{enumerate}
  \item Show that $N$ is the midpoint of $[AC]$.
  \item Explain why $(MN)$ is \omterm{def:g6:lines:perp}{perpendicular} to $(AC)$.
\end{enumerate}
\end{exercise}

\begin{exercise}[$\star\star\star$]\label{exo:g8:midpoints:9}
Let $ABCD$ be \emph{any} \omterm{def:g4:geometry:polygon}{quadrilateral}, and $I$, $J$, $K$, $L$ the
midpoints of $[AB]$, $[BC]$, $[CD]$, $[DA]$ (the conjecture of
\cref{exo:g7:central:11}).
\begin{enumerate}
  \item Apply the midpoint theorem in the triangle $ABC$ to the
        \omterm{def:g6:lines:objects}{segment} $[IJ]$, and in the triangle $ACD$ to the \omterm{def:g6:lines:objects}{segment}
        $[LK]$: compare each to the diagonal $[AC]$.
  \item Conclude that $IJKL$ is always a \omterm{def:g7:central:parallelogram}{parallelogram}.
\end{enumerate}
\end{exercise}
